\paragraph{}The main goal of this lab was to better understand how three-phase transformers are configured and how transformers are used to supply power to residential homes. In the first part of the lab, we were able to see how the Delta-Wye configuration affected the line and phase voltages and currents on both the primary and secondary sides of the transformer bank.
In the second part of the lab, we were able to see how different loading conditions affected a residential transformer. When both 120 V loads were connected, very little current was measured through the neutral. Once one of the loads was disconnected, the current through the neutral increased due to the imbalance between the two loads. We also saw a decrease in the secondary voltages as more loads were connected. Even with all of the loads disconnected, the transformer was still consuming a small amount of power due to internal losses.
Overall, this lab gave us a better understanding of the relationships between voltage, current, and power in different transformer configurations. It also helped demonstrate the importance of balancing loads in residential power systems and how loading can affect the performance of a transformer.

